% ECONOMICS_PROBLEM: {"id": "J52-253", "title": "Collective-bargaining redesign", "jel_code": "J52", "source_id": "OP-0354"}
% !TeX program = xelatex
\documentclass[11pt,a4paper]{article}
\usepackage[margin=23mm]{geometry}
\usepackage{fontspec}
\setmainfont{Latin Modern Roman}
\usepackage{amsmath,amssymb,mathtools}
\usepackage{xcolor,enumitem,booktabs,longtable,array,xurl,hyperref}
\definecolor{muted}{HTML}{53636C}
\hypersetup{colorlinks=true,urlcolor=blue,linkcolor=blue}
\setlength{\parindent}{0pt}
\setlength{\parskip}{5pt}
\newcommand{\R}{\mathbb R}
\newcommand{\E}{\mathbb E}
\newcommand{\N}{\mathbb N}
\newcommand{\ind}{\mathbf 1}
\newcommand{\Var}{\operatorname{Var}}
\newcommand{\Cov}{\operatorname{Cov}}
\newcommand{\supp}{\operatorname{supp}}
\DeclareMathOperator*{\argmin}{arg\,min}
\DeclareMathOperator*{\argmax}{arg\,max}
\newcommand{\entrymeta}[3]{{\sffamily\small\color{muted}\textbf{\texttt{#1}}\quad|\quad #2\par #3\par}}
\newcommand{\Problemref}[1]{\texttt{#1}}
\newcommand{\JELref}[2]{\texttt{#2} (JEL \texttt{#1})}
\begin{document}
\section*{Collective-bargaining redesign}
\textbf{Problem ID:} \texttt{J52-253}\quad \textbf{JEL:} \texttt{J52}\par
% BEGIN ECONOMICS STATEMENT
\label{problem:OP-0354}
{\sffamily\small\color{muted}Original subject group: Labor markets and work\par}
\label{definitions:problem:30007636}
\textbf{Audit classification:} Published research agenda. Published metadata and full text checked. Political-policy and coverage questions are explicit; optimization over bargaining systems is editorial.
\subsection*{Definitions and precise research target}
\textbf{Complete editorial problem.} Consider a specified economy with workers, firms, and a finite set \(\mathcal B\) of bargaining arrangements. An arrangement \(b\) defines bargaining coverage, representation rules, agreement extension, contract enforcement, and negotiation level. Let \(W_T(b)\), \(L_T(b)\), \(K_T(b)\), and \(R_T(b)\) denote equilibrium compensation, employment, investment, and a preregistered representation index at horizon \(T\). Let \(U_i(b)\) be worker \(i\)'s lifetime utility, measured with an explicit wage--amenity preference model, and let \(\omega_i\ge0\) be declared welfare weights summing to one. Define \(S(b)=E[\sum_i\omega_iU_i(b)]\). The design task is to characterize arrangements maximizing \(S(b)\) subject to specified feasibility, fiscal, and employment constraints. Estimate effects separately for represented workers, uncovered workers, entrants, and consumers. A credible analysis must allow bargaining to change firm entry, productivity, wage-setting, political influence, and who becomes represented. State the variation that distinguishes these mechanisms and the equilibrium closure used in counterfactuals. Existing evidence about a particular union's wage premium cannot identify this entire policy ordering. The optimum is conditional on stated objectives and institutional context, rather than a universally best bargaining system.

\textbf{Definitions and admissible scope.} Let the finite arrangement menu specify a separate coverage indicator, membership indicator, negotiation level, extension rule, and dispute procedure. Worker utilities \(U_i(b)\) are cardinally normalized using stated consumption, amenity, and risk preferences; wages alone are not utilities. Define normalized weights \(\omega_i\ge0\), worker-hours \(L_T(b)\), real investment \(K_T(b)\), and representation \(R_T(b)\), for example the share of eligible workers able to vote for negotiators. Feasibility requires stated inequalities \(L_T(b)\ge\underline L\) and \(C(b)\le B\). The optimum of a known finite menu is elementary. The substantive task is to identify equilibrium outcomes and welfare rankings, permitting bargaining to affect uncovered workers, consumers, entry, and policy. Define the equilibrium or reduced-form intervention law before interpreting a collective-bargaining counterfactual. The population, objective weights, arrangements, feasibility constraints, utility functions, real resource costs, and treatment of uncertainty are all fixed inputs. Take \(\mathcal B\) finite and \(E|U_i(b)|<\infty\) for every \(i,b\). The requested output is the identified set of feasible welfare rankings and maximizers when equilibrium response laws are not known, not merely enumeration after all utilities have been supplied. The displayed \(S\) is explicitly a worker-weighted criterion; a whole-economy variant must add separately weighted consumer, owner, and fiscal payoffs and avoid counting transfers as additional resources.
\subsection*{Variants and assumption changes}
\begin{enumerate}
\item Coverage versus membership (source-motivated): vary agreement coverage while holding membership rules fixed, then vary membership conditional on coverage. Estimate separate wage and political-participation effects; overlapping interventions require a factorial design or additional restrictions.
\item Political feedback (source-motivated extension): let arrangement \(b\) change subsequent labor legislation through an explicit voting or lobbying model. Compare the endogenous-policy welfare ranking with a ranking that holds legislation fixed; neither ranking is universally best.
\end{enumerate}
\subsection*{References and source audit}
\textbf{Support and access.} Published metadata and full text checked. Political-policy and coverage questions are explicit; optimization over bargaining systems is editorial. \textbf{Locator:} Section 7, pp. 389--390, policy and coverage-versus-membership bullets.
\begin{enumerate}\raggedright
\item\hypertarget{definitions:source:30007636:1}{} Ethan Kaplan; Suresh Naidu. 2025. \emph{Between Government and Market: The Political Economics of Labor Unions}. Section 7, Conclusion and Future Research Questions, printed pp. 389–390, bullets on direct policy effects and coverage versus membership.
\url{https://econweb.umd.edu/~kaplan/PE_of_Labor.pdf}.
DOI: \url{https://doi.org/10.1146/annurev-economics-051520-025251}.
\end{enumerate}

\subsection*{Common conventions: Source mathematical conventions}

These conventions apply to each entry unless explicitly replaced.

\textbf{Models and identification.} A primitive model \(m\) specifies all probability laws, feasible choices, information, equilibrium and measurement rules used by its target. The admissible class \(\mathcal M\) is fixed independently of outcomes. If \(P_O(m)\) is its observable probability law and \(\tau(m)\) the target, its identified set at observable law \(P\) is
\[
 \mathcal I_\tau(P)=\{\tau(m):m\in\mathcal M,\ P_O(m)=P\}.
\]
Point identification means that this set is a singleton. An empty set signals incompatibility of data and assumptions. Coordinate bounds need not describe the joint set. Bounds are sharp only if every retained target is attainable or arbitrarily approachable by compatible models. Finite bounds require explicit restrictions on unobserved quantities; unrestricted confounding, hidden wealth, extrapolation or arbitrary equilibrium selection can yield uninformative sets.

\textbf{Causal interventions.} A potential outcome \(Y_i(p)\) is the outcome under a complete policy regime \(p\), including timing, financing, information and market spillovers. The target population and units are fixed before treatment. With interference, use the complete assignment vector or a justified exposure mapping; individual no-interference notation alone cannot identify an economy-wide intervention. A difference of mean potential outcomes does not require the cross-world joint distribution. Conditioning on post-treatment migration, survival, enrollment or employment changes the estimand and needs separate justification.

\textbf{Statistical statements.} An estimator, confidence set, forecast or test specifies a sampling law, model class and loss/coverage criterion. Uniform \((1-\alpha)\)-coverage means \(\inf_{m\in\mathcal M}\Pr_m\{\tau(m)\in C_n\}\geq1-\alpha\), or an explicitly asymptotic version. The whole parameter space trivially has coverage, so informativeness requires a stated width, risk or power objective. A forecasting improvement is relative to a named benchmark, information set and evaluation population.

\textbf{Equilibrium and optimization.} An equilibrium comprises feasible plans, optimality, beliefs where needed, budget constraints and all market/resource clearing. The primitives or a stated benchmark define each correspondence \(\mathcal E(p;m)\). Nonemptiness is assumed or proved, never inferred just from compact policy sets. A selected-equilibrium target uses a declared selection rule; a robust target quantifies over all permitted equilibria. Use suprema or infima unless attainment is established. An \(\varepsilon\)-optimal policy is feasible and within \(\varepsilon\) of the finite optimum value. Report an empty feasible set separately.

\textbf{Welfare and accounting.} Normative utility, interpersonal normalization, social weights, numeraire, discounting and the use of public receipts are primitives. Behavioral utility need not coincide with the welfare criterion. Household welfare includes ownership income and rents once; adding profits again to compensating variation generally double-counts them. A monetary equivalent is defined by an explicit transfer and price convention, with monotonicity and range conditions. Average effects, mechanism contributions, transfers and welfare losses are different targets. Infinite sums require convergence and dynamic feasible plans require terminal or no-Ponzi conditions.

\textbf{Source versions.} A working-paper date, revision date and journal publication year are distinguished. A DOI for a journal version is not automatically the DOI of an earlier manuscript. Locators refer to the stated version and printed pages where specified. A metadata-only check does not verify a claim about a full-text conclusion. Every entry's source subsection narrows the provenance claim accordingly.
% END ECONOMICS STATEMENT
\end{document}
